\documentclass[10pt,a4paper]{article}
\usepackage[T1]{fontenc}
\usepackage{lmodern}
\usepackage{amsmath,amssymb,amsthm,mathtools}
\usepackage[margin=24mm]{geometry}
\usepackage{microtype}
\usepackage{booktabs,array,longtable}
\usepackage{xcolor}
\usepackage{hyperref}
\usepackage{fancyhdr}
\hypersetup{colorlinks=true,linkcolor=black,citecolor=black,urlcolor=blue,pdfauthor={Matthew J. Colbrook}}
\pagestyle{fancy}
\fancyhf{}
\fancyhead[L]{\small Resolution / MI-19}
\fancyhead[R]{\small 11 September 2026}
\fancyfoot[C]{\thepage}
\setlength{\headheight}{14pt}
\setlength{\parskip}{4pt}
\setlength{\parindent}{0pt}
\newtheorem{theorem}{Theorem}[section]
\newtheorem{proposition}[theorem]{Proposition}
\newtheorem{lemma}[theorem]{Lemma}
\newtheorem{corollary}[theorem]{Corollary}
\theoremstyle{remark}
\newtheorem{remark}[theorem]{Remark}
\newcommand{\M}{\mathbb M}
\newcommand{\C}{\mathbb C}
\newcommand{\R}{\mathbb R}
\newcommand{\tr}{\operatorname{tr}}
\newcommand{\diag}{\operatorname{diag}}
\newcommand{\spec}{\operatorname{spec}}
\newcommand{\rank}{\operatorname{rank}}
\newcommand{\abs}[1]{\lvert #1\rvert}
\newcommand{\norm}[1]{\lVert #1\rVert}
\newcommand{\opnorm}[1]{\lVert #1\rVert_{\infty}}
\newcommand{\schatten}[2]{\lVert #1\rVert_{#2}}
\newcommand{\symmod}{\mathcal S}
\newcommand{\maxmod}{\mathcal M}
\newcommand{\draftnotice}{\begin{center}\small\textit{Proposed mathematical resolution. Complete proof supplied; not submitted or peer reviewed.}\end{center}}

\usepackage{xurl}
\setlength{\emergencystretch}{3em}
\allowdisplaybreaks[1]

\title{MI-19: Resolution}
\author{Matthew J. Colbrook\\
\small Department of Applied Mathematics and Theoretical Physics\\
\small University of Cambridge, Cambridge, United Kingdom\\
\small\href{mailto:m.colbrook@damtp.cam.ac.uk}{m.colbrook@damtp.cam.ac.uk}}
\date{11 September 2026}
\begin{document}
\maketitle
\textbf{Repository status: Solved.}
The complete argument received an independent Codex-agent PASS review
against the exact catalog target.
The original draft was AI-assisted; the reviewed proof below is unchanged.
This is independent agent verification, not external human peer review or
formal proof certification. \href{https://github.com/ajt60gaibb/OpenProblemsInNLA/blob/main/references/colbrook-matrix-2026-09-11/verification/reviews/MI-19-review.md}{Detailed review and scope}.
\par\medskip
% BEGIN REVIEWED PROOF
\section{MI-19: a rank-two counterexample to the subset inequality for the \texorpdfstring{$q$}{q}-permanent}
\label{sec:mi19}
For $A=(a_{ij})\in\M_n$ and $q\in\R$, use the inversion convention
\[
 \operatorname{per}_q(A)=\sum_{\sigma\in S_n}q^{\ell(\sigma)}\prod_{i=1}^n a_{i,\sigma(i)},\qquad
 \ell(\sigma)=\#\{(i,j):i<j,\ \sigma(i)>\sigma(j)\}.
\]
For a subset $S\subseteq\{1,\ldots,n\}$, define the restricted sum
\[
 R_{q,S}(A)=\sum_{\substack{\sigma\in S_n\\\sigma(S)=S}}
 q^{\ell(\sigma)}\prod_{i=1}^n a_{i,\sigma(i)}.
\]
The target is $\operatorname{per}_q(A)\geq R_{q,S}(A)$ for Hermitian positive semidefinite $A$ and $0\leq q\leq1$. This is the subset formulation recorded in \cite[Conjecture 3, equation (5.1)]{Fonseca2018}.

\begin{theorem}
The subset inequality is false for a real positive semidefinite matrix of order four and rank two. A counterexample is
\[
 A=\begin{pmatrix}
 170&1&167&170\\
 1&1&-2&1\\
 167&-2&173&167\\
 170&1&167&170
 \end{pmatrix},\qquad q=\frac78,\qquad S=\{2\}.
\]
\end{theorem}
\begin{proof}
Positive semidefiniteness and rank two follow from
\[
 A=X^TX,\qquad X=\begin{pmatrix}13&0&13&13\\1&1&-2&1\end{pmatrix}.
\]
Expansion over the $24$ permutations gives
\begin{align*}
 \operatorname{per}_q(A)
 &=115600q^6+4886140q^5+9568969q^4+4712758q^3\\
 &\hspace{1.5em}-199231q^2+4886140q+4999700.
\end{align*}
The six permutations fixing the second index give
\[
 R_{q,\{2\}}(A)=4999700q^5+9482260q^4+4741130q^3+4741130q+4999700.
\]
Consequently the difference is the small polynomial
\begin{align*}
 \operatorname{per}_q(A)-R_{q,\{2\}}(A)
 &=q(q+1)\bigl(115600q^4-229160q^3+315869q^2\\
 &\hspace{8em}-344241q+145010\bigr).
\end{align*}
At $q=7/8$ this equals
\[
 -\frac{3235575}{16384}<0.
\]
All calculations are integer polynomial arithmetic. The accompanying verifier independently enumerates the permutations with the stated inversion convention.
\end{proof}

\begin{remark}[Positive definite variants and scope]
The strict negative gap persists for $A+\varepsilon I_4$ for every sufficiently small positive $\varepsilon$, by continuity of the two polynomial expressions. Thus the obstruction is not dependent on allowing a singular input. No counterexample to the separate monotonicity problem MI-18 is asserted.
\end{remark}

% END REVIEWED PROOF
\begin{thebibliography}{1}

\bibitem{Fonseca2018}
Carlos~M. da~Fonseca.
\newblock The $\mu$-permanent revisited.
\newblock arXiv:1804.02231, 2018.
\newblock Conjecture 3, equation (5.1).
\newblock URL: \url{https://arxiv.org/abs/1804.02231}.

\end{thebibliography}

\end{document}
